Douglas Rodrigues Barros (2026) e Gabriel Santos (2026) leem a oposição entre monoteísmo e politeísmo em Nego Bispo como uma explicação religiosa da escravização e contestam o peso que ele atribuiria à religião. Há fundamento para essa leitura: em conversa com Renato Noguera, Bispo (2022) distingue a necessidade colonial de mão de obra da pergunta sobre quem determinou quais povos poderiam ser escravizados e, para responder à segunda, remete à bula papal de 1455. Mas, quando diz a Gabriel Galo que um povo de cosmologia politeísta \enquote{pode ser até cristão, mas não é monoteísta}, Bispo (2020) abre a questão que me interessa: além de dizer a quem se presta culto, o que essas palavras dizem sobre quem pode agir junto - os vivos, os ancestrais e os demais seres com os quais se compõe um mundo?